\section{Formalization Map}
\label{appendix:formalization-map}

This appendix records which statements of the paper are machine-checked in
the Lean~4 development \texttt{FastPoly/\allowbreak } distributed with the source, under
which hypotheses, and where the formalization stops short of the paper.  The
abstract's sentence ``Everything is checked in Lean'' refers to the complete
upper-bound construction stated immediately before it: the decoder,
multiplication count, addition bounds, and logarithmic-height arrangement.
Those claims are joined by the paper-level capstone in
\texttt{PaperMain.lean}.  The sentence does not refer to independent auxiliary
results recorded later in the paper, whose coverage is reported separately
below.
The
Lean statements are made over an abstract commutative ring $R$ and a
commutative $R$-algebra $A$ (Lean: \texttt{[CommRing R] [CommRing A]
[Algebra R A]}); the only characteristic-type assumptions are explicit
\texttt{IsUnit} hypotheses on integers, discussed in
\Cref{sec:formalization-map:units}.  Decoding is expressed as membership in
a generated subalgebra: ``$\theta$ is recovered from $P$'' is
$\theta\in K\sqcup\operatorname{adjoin}_R\{[x^i]P\}$ for the ambient known
algebra $K$, and every recovery proof exhibits the pivot, peel, division, or
block inverse it uses (the repository forbids search and decision
procedures on nontrivial goals).  In the tables, \emph{full} means the
displayed paper statement is proved as stated (up to the deviations noted in
the hypotheses column), \emph{partial} means a strictly weaker or
differently packaged statement is proved, and \emph{not formalized} means
no Lean counterpart exists.  File names are relative to \texttt{FastPoly/\allowbreak }.

\subsection{Main results}
\label{sec:formalization-map:main}

\begingroup
\footnotesize
\setlength{\tabcolsep}{3pt}
\par\noindent\begin{tabular}{p{0.95in}p{2.05in}p{1.75in}p{0.55in}}
\toprule
Paper & Lean (file) & Hypotheses in Lean & Status\\
\midrule
\end{tabular}\par\nopagebreak
\noindent\begin{tabular}{p{0.95in}p{2.05in}p{1.75in}p{0.55in}}
\raggedright \Cref{thm:main} &\raggedright
\texttt{MainArrangementsChecked.of\_\allowbreak charZero},
\texttt{MainArrangementsChecked.of\_\allowbreak charP}
(\texttt{PaperMain.lean});
\texttt{decodedAdditionPolynomial\_\allowbreak exists}
(\texttt{Cost/\allowbreak Additions/\allowbreak DecodedPolynomial.lean});
\texttt{monic\_\allowbreak coefficient\_\allowbreak map\_\allowbreak bijective}
(\texttt{Main.lean});
\texttt{polynomial\_\allowbreak height} (\texttt{HeightFinal.lean});
\texttt{MvPolynomial.algEquivOfDecodable} (\texttt{Automorphism.lean}) &
\raggedright Characteristic zero or $p>n$.  The paper-level capstone constructs
both advertised arrangements at once on one common monic degree-$n$ polynomial:
a fixed program carrying the decoder, multiplication count, and logarithmic height;
and a fixed addition-optimized program carrying its semantics, multiplication bound,
exact literal addition count, and both addition inequalities.  The two scheduling
objectives may use separate fixed syntax trees, but both are proved to compute the
same polynomial;
the coefficient preprocessing map is the polynomial automorphism of
\Cref{lem:polynomial-left-inverse-automorphism}. &
full\\
\end{tabular}\par
\noindent\begin{tabular}{p{0.95in}p{2.05in}p{1.75in}p{0.55in}}
\raggedright \Cref{thm:odd-realizable-pairs} &\raggedright \texttt{odd\_\allowbreak realizable\_\allowbreak pairs}, \texttt{odd\_\allowbreak realizable\_\allowbreak pairs'}
(\texttt{Main.lean}) &
\raggedright \texttt{[Nontrivial A]}; $n$ odd, $3\le n$, $n\ne7$;
\texttt{$\forall$ i, 1 $\le$ i $\to$ i $\le$ n $\to$ IsUnit ((i : $\mathbb Z$) : R)}; any parameter block
\texttt{$\theta$ : $\mathbb N$ $\to$ A}.  Conclusion: \texttt{CompatiblePair $\bot$ T$_1$ T$_2$ (n-1) G};
$H_2$ (and $H_4$ when $n\ge5$) monic of degree $2$ ($4$); every subalgebra
containing the coefficients of $x T_1+T_2$ contains $\theta_0,\dots,\theta_{n-1}$
and the coefficients of $H_2,H_4$; and a program
\texttt{prog : Cost.\allowbreak JointPairProgram R ((n-1)/\allowbreak 2)} with
\texttt{prog.\allowbreak RealizesAt $\theta$ T$_1$ T$_2$ H$_2$ H$_4$} and
\texttt{prog.\allowbreak HeightBounded (2 * Nat.clog 2 n + 3)}. &
full\\
\end{tabular}\par
\noindent\begin{tabular}{p{0.95in}p{2.05in}p{1.75in}p{0.55in}}
\raggedright \Cref{thm:odd-realizable-pairs}, fixed program before keys &\raggedright \texttt{odd\_\allowbreak realizable\_\allowbreak pairs\_\allowbreak free},
\texttt{odd\_\allowbreak realizable\_\allowbreak pairs\_\allowbreak uniform\_\allowbreak family} (\texttt{Main.lean}) &
\raggedright Instantiation at $A=R'[\alpha_0,\dots,\alpha_{n-1}]$
(\texttt{MvPolynomial (Fin n) R'}); the same program syntax realizes the
specialized pair and powers for every key vector over every $R'$-algebra
(\texttt{RealizesFiniteFamily}). &
full\\
\end{tabular}\par
\noindent\begin{tabular}{p{0.95in}p{2.05in}p{1.75in}p{0.55in}}
\raggedright \Cref{thm:construction-count}, count clause &\raggedright \texttt{Cost.pair\_\allowbreak construction\_\allowbreak count},
\texttt{Cost.construction\_\allowbreak count} (\texttt{Cost/\allowbreak Final.lean});
\texttt{Cost.\allowbreak SepticProgram.program} (\texttt{Cost/\allowbreak SepticProgram.lean});
\texttt{Cost.\allowbreak PolynomialProgram.evenLift},
\texttt{evenLift\_\allowbreak realizesAt} (\texttt{Cost/\allowbreak PolynomialProgram.lean}) &
\raggedright Recurrences \texttt{PairCost}/\texttt{PolynomialCost} over $\mathbb N$ (no ring):
$\exists c$, \texttt{PolynomialCost n c} with $c=\lfloor n/2\rfloor+1$ for
$n\ge3$, $c=0,1$ for $n=1,2$.  The literal program with $(n-1)/2$ products is
the \texttt{JointPairProgram} conjunct of the master; the septic program has
four products; the even lift adds one product and one addition. &
full\\
\end{tabular}\par
\noindent\begin{tabular}{p{0.95in}p{2.05in}p{1.75in}p{0.55in}}
\raggedright \Cref{thm:construction-count}, coverage clause;
\Cref{cor:all-odd-decodable} &\raggedright \texttt{monic\_\allowbreak coefficient\_\allowbreak map\_\allowbreak bijective},
\texttt{even\_\allowbreak lift\_\allowbreak bijective},
\texttt{odd\_\allowbreak coefficient\_\allowbreak map\_\allowbreak bijective} (\texttt{Main.lean}) &
\raggedright \texttt{[IsNoetherianRing R'] [Nontrivial R']}; $1,\dots,n$ units; every
$n\ge1$ (affine, quadratic, septic, odd master, even lift).  Conclusion: a
monic degree-$n$ $P$ over the free coordinate algebra whose coefficient
substitution \texttt{MvPolynomial.aeval (fun i => P.coeff i)} is bijective. &
full\\
\end{tabular}\par
\noindent\begin{tabular}{p{0.95in}p{2.05in}p{1.75in}p{0.55in}}
\raggedright \Cref{thm:construction-height} &\raggedright
\texttt{polynomial\_\allowbreak height},
\texttt{odd\_\allowbreak polynomial\_\allowbreak height},
\texttt{septic\_\allowbreak polynomial\_\allowbreak height},
\texttt{RealizedPolynomial.\allowbreak evenLift},
\texttt{polynomial\_\allowbreak height\_\allowbreak of\_\allowbreak charZero},
\texttt{polynomial\_\allowbreak height\_\allowbreak of\_\allowbreak charP}
(\texttt{HeightFinal.lean}); height conjunct \texttt{HeightBounded} of
\texttt{odd\_\allowbreak realizable\_\allowbreak pairs} (\texttt{Main.lean});
\texttt{Circuit.multDepth} (\texttt{Height/\allowbreak Depth.lean}); component
ledgers in \texttt{Height/\allowbreak PeeledCircuit.lean} and
\texttt{Height/\allowbreak TCircuitDepth.lean} &
\raggedright \texttt{[Nontrivial A]}; $3\le n$; $1,\dots,n$ units of $R$ (for the
\texttt{\_of\_charZero}/\texttt{\_of\_charP} forms: a field of characteristic
zero, resp.\ $p>n$).  Conclusion
\texttt{RealizedPolynomial R $\theta$ n (n/2+1) (2 * Nat.clog 2 n + 4 + (n+1) \% 2)}:
one fixed complete-polynomial program is monic of degree $n$, decodes all $n$
parameters, uses exactly $\lfloor n/2\rfloor+1$ products, and has output
\texttt{multDepth} (inputs at depth $0$) at most $2\lceil\log_2 n\rceil+4$ for
odd $n$ and $2\lceil\log_2 n\rceil+5$ for even $n$.  Assembled from the
master's pair conjunct ($T_1,T_2$ at depth $\le2\lceil\log_2 n\rceil+3$, $H_2$
at $\le1$, $H_4$ at $\le2$) by the final product $xT_1+T_2$, the direct septic
($n=7$, lifted for $n=8$), and the even lift $xQ+c_0$.  The stronger
\texttt{DecodedAdditionPolynomial} capstone pairs this height witness with the
addition-certified arrangement on the same semantic polynomial. &
full (count, decoder, height)\\
\end{tabular}\par
\noindent\begin{tabular}{p{0.95in}p{2.05in}p{1.75in}p{0.55in}}
\raggedright \Cref{thm:construction-addition-count},
\Cref{thm:construction-addition-asymptotic} &\raggedright \texttt{Cost.construction\_\allowbreak additions\_\allowbreak checked}
(\texttt{Cost/\allowbreak Additions/\allowbreak Realization.lean}); ledger
\texttt{Cost.construction\_\allowbreak addition\_\allowbreak count}
(\texttt{Cost/\allowbreak Additions/\allowbreak Final.lean}) &
\raggedright For every $n\ge1$ over an $n$-admissible ring and every parameter
environment: one fixed program \texttt{program : PolynomialProgram R mult}
with \texttt{mult}\,$\le\lfloor n/2\rfloor+1$, \texttt{program.RealizesAt}
its degree-$n$ polynomial, and its literal gate count
\texttt{program.additions} satisfying both $\le2n$ and
$4\cdot{}\le5n+24\lceil\log_2n\rceil^2+4$.  Assembled from the
addition-certified base pairs, the recursive $8k+3$/$8k+7$ pair steps, the
certified odd-gadget dispatch, and the odd/even combinators; the program is
the optimized schedule of \Cref{sec:addition-count}, not the master's. &
full\\
\bottomrule
\end{tabular}\par
\endgroup

\subsection{Decoder calculus\texorpdfstring{ (\Cref{appendix:decoder-calculus})}{}}
\label{sec:formalization-map:calculus}

\begingroup
\footnotesize
\setlength{\tabcolsep}{3pt}
\par\noindent\begin{tabular}{p{0.95in}p{2.05in}p{1.75in}p{0.55in}}
\toprule
Paper & Lean (file) & Hypotheses in Lean & Status\\
\midrule
\end{tabular}\par\nopagebreak
\noindent\begin{tabular}{p{0.95in}p{2.05in}p{1.75in}p{0.55in}}
\raggedright \Cref{def:compatible-pair}; visible algebra &\raggedright \texttt{CompatiblePair}, \texttt{Vis}, \texttt{CausalPair}
(\texttt{Recover/\allowbreak Context.lean}) &
\raggedright \texttt{[CommRing R] [CommRing A] [Algebra R A]}; relative to a known
subalgebra $K$ and a window $G$. &
full\\
\end{tabular}\par
\noindent\begin{tabular}{p{0.95in}p{2.05in}p{1.75in}p{0.55in}}
\raggedright scalar pivots \eqref{eq:filtered-pivot} &\raggedright \texttt{mem\_\allowbreak obsAlg\_\allowbreak of\_\allowbreak scalarCert} (\texttt{Recover/\allowbreak Filtered.lean}) &
\raggedright Each pivot slope \texttt{lam j : R} is a unit; well-founded index order. &
full\\
\end{tabular}\par
\noindent\begin{tabular}{p{0.95in}p{2.05in}p{1.75in}p{0.55in}}
\raggedright block pivots with explicit inverse &\raggedright \texttt{mem\_\allowbreak of\_\allowbreak blockCert} (\texttt{Recover/\allowbreak Filtered.lean});
\texttt{mem\_\allowbreak of\_\allowbreak known\_\allowbreak blockCert\_\allowbreak of\_\allowbreak det} (\texttt{Recover/\allowbreak KnownBlock.lean}) &
\raggedright Constant matrix $M$ over $R$ with a supplied $N$, $NM=1$; or entries in the
known algebra with $\det M$ a unit of $R$ (adjugate inverse). &
full\\
\end{tabular}\par
\noindent\begin{tabular}{p{0.95in}p{2.05in}p{1.75in}p{0.55in}}
\raggedright \Cref{def:coefficient-triangular}, \Cref{lem:triangular-shift},
\Cref{lem:triangular-implies-compatible} &\raggedright \texttt{CoeffTriangular}, \texttt{CoeffTriangular.shift\_\allowbreak pivot},
\texttt{CoeffTriangular.toCompatiblePair} (\texttt{Recover/\allowbreak Triangular.lean}) &
\raggedright Unit pivots \texttt{lam j} for $j<d$; the compatibility half is proved. &
full\\
\end{tabular}\par
\noindent\begin{tabular}{p{0.95in}p{2.05in}p{1.75in}p{0.55in}}
\raggedright \Cref{lem:triangular-block-concatenation} &\raggedright \texttt{adjoin\_\allowbreak Ico\_\allowbreak glue}, \texttt{side\_\allowbreak data\_\allowbreak absorb}
(\texttt{Recover/\allowbreak Triangular.lean}) &
\raggedright Proved as the two glue lemmas used at every concatenation site. &
full\\
\end{tabular}\par
\noindent\begin{tabular}{p{0.95in}p{2.05in}p{1.75in}p{0.55in}}
\raggedright \Cref{lem:monic-cauchy-transport} &\raggedright \texttt{coeff\_\allowbreak mul\_\allowbreak monic} (\texttt{Polynomial/\allowbreak TopWindow.lean}) &
\raggedright \texttt{[CommRing A]}; $q$ monic.  No unit hypothesis. &
full\\
\end{tabular}\par
\noindent\begin{tabular}{p{0.95in}p{2.05in}p{1.75in}p{0.55in}}
\raggedright \Cref{lem:peel-monic-factor} &\raggedright \texttt{coeff\_\allowbreak mem\_\allowbreak of\_\allowbreak mul\_\allowbreak monic\_\allowbreak add} (\texttt{Polynomial/\allowbreak PeelMonic.lean}) &
\raggedright $M$ monic with coefficients in $K$.  No unit hypothesis. &
full\\
\end{tabular}\par
\noindent\begin{tabular}{p{0.95in}p{2.05in}p{1.75in}p{0.55in}}
\raggedright \Cref{lem:monic-from-power} &\raggedright \texttt{coeff\_\allowbreak mem\_\allowbreak of\_\allowbreak pow\_\allowbreak add} (\texttt{Polynomial/\allowbreak MonicFromPower.lean}) &
\raggedright \texttt{IsUnit (m : R)} for the root index $m$. &
full\\
\end{tabular}\par
\noindent\begin{tabular}{p{0.95in}p{2.05in}p{1.75in}p{0.55in}}
\raggedright \Cref{lem:monic-division} &\raggedright \texttt{coeff\_\allowbreak quot\_\allowbreak mem}, \texttt{coeff\_\allowbreak rem\_\allowbreak mem}
(\texttt{Polynomial/\allowbreak MonicDivision.lean}) &
\raggedright Divisor monic with coefficients in $K$.  No unit hypothesis. &
full\\
\end{tabular}\par
\noindent\begin{tabular}{p{0.95in}p{2.05in}p{1.75in}p{0.55in}}
\raggedright \Cref{lem:square-gadget}, \Cref{lem:square-gadget-boundary} &\raggedright \texttt{square\_\allowbreak gadget\_\allowbreak mem} (\texttt{Polynomial/\allowbreak SquareGadget.lean});
\texttt{coeff\_\allowbreak mem\_\allowbreak of\_\allowbreak square\_\allowbreak gadget\_\allowbreak relative}
(\texttt{Polynomial/\allowbreak CausalShell.lean}) &
\raggedright \texttt{IsUnit (2 : R)}. &
full\\
\end{tabular}\par
\noindent\begin{tabular}{p{0.95in}p{2.05in}p{1.75in}p{0.55in}}
\raggedright \Cref{lem:scalar-shift-square} &\raggedright \texttt{scalar\_\allowbreak shift\_\allowbreak mem} (\texttt{Polynomial/\allowbreak ScalarShift.lean}) &
\raggedright \texttt{IsUnit lam} and \texttt{IsUnit (2 : R)}. &
full\\
\end{tabular}\par
\noindent\begin{tabular}{p{0.95in}p{2.05in}p{1.75in}p{0.55in}}
\raggedright \Cref{lem:x-alpha-extraction} &\raggedright \texttt{x\_\allowbreak alpha\_\allowbreak mem} (\texttt{Recover/\allowbreak XAlpha.lean}) &
\raggedright A compatible pair with window $G\subseteq[0,n)$.  No unit hypothesis. &
full\\
\end{tabular}\par
\noindent\begin{tabular}{p{0.95in}p{2.05in}p{1.75in}p{0.55in}}
\raggedright \Cref{lem:discharge-side-information},
\Cref{lem:extractable-via-derivable} &\raggedright \texttt{discharge\_\allowbreak side\_\allowbreak information},
\texttt{extractable\_\allowbreak via\_\allowbreak derivable} (\texttt{Recover/\allowbreak Context.lean}) &
\raggedright As stated. &
full\\
\end{tabular}\par
\noindent\begin{tabular}{p{0.95in}p{2.05in}p{1.75in}p{0.55in}}
\raggedright \Cref{lem:polynomial-left-inverse-automorphism} &\raggedright \texttt{MvPolynomial.algEquivOfDecodable} (\texttt{Automorphism.lean});
\texttt{coefficientAlgEquiv\allowbreak OfMonicDecodable} (\texttt{Instantiation.lean}) &
\raggedright \texttt{[IsNoetherianRing R] [Finite $\sigma$]}; decodability of the variables from
the coefficient family gives a two-sided $R$-algebra automorphism. &
full\\
\end{tabular}\par
\noindent\begin{tabular}{p{0.95in}p{2.05in}p{1.75in}p{0.55in}}
\raggedright \Cref{def:joint-realization} &\raggedright \texttt{Cost.\allowbreak JointPairProgram}, \texttt{JointPairProgram.\allowbreak RealizesAt}
(\texttt{Cost/\allowbreak PolynomialCircuit.lean});
\texttt{Cost.\allowbreak MultiplicationProgram} (\texttt{Cost/\allowbreak MultiplicationProgram.lean}) &
\raggedright A straight-line program over $R$ with a stored multiplication count; four
outputs $(T_1,T_2,H_2,H_4)$. &
full\\
\bottomrule
\end{tabular}\par
\endgroup

\subsection{Construction layers\texorpdfstring{ (\Cref{appendix:constructions})}{}}
\label{sec:formalization-map:construction}

\begingroup
\footnotesize
\setlength{\tabcolsep}{3pt}
\par\noindent\begin{tabular}{p{0.95in}p{2.05in}p{1.75in}p{0.55in}}
\toprule
Paper & Lean (file) & Hypotheses in Lean & Status\\
\midrule
\end{tabular}\par\nopagebreak
\noindent\begin{tabular}{p{0.95in}p{2.05in}p{1.75in}p{0.55in}}
\raggedright \Cref{alg:constr-fill}, \Cref{lem:fill-correctness} &\raggedright \texttt{fill\_\allowbreak correct} (\texttt{Section4/\allowbreak FillRec.lean}) &
\raggedright \texttt{[Nontrivial A]}; per-level certificates \texttt{GoodLevel}; known
monic quadratic $H_1$. &
full\\
\end{tabular}\par
\noindent\begin{tabular}{p{0.95in}p{2.05in}p{1.75in}p{0.55in}}
\raggedright \Cref{alg:constr-known-2n-1}, \Cref{lem:Q-unitriangular} &\raggedright \texttt{peel}, \texttt{peel\_\allowbreak correct} (\texttt{Section4/\allowbreak Peeled.lean});
\texttt{peel\_\allowbreak unitriangular} (\texttt{Section4/\allowbreak PeeledCert.lean}) &
\raggedright \texttt{[Nontrivial A]}; a tower of known monic powers of degrees $2^i$. &
full\\
\end{tabular}\par
\noindent\begin{tabular}{p{0.95in}p{2.05in}p{1.75in}p{0.55in}}
\raggedright \Cref{lem:peeled-Q-decodable}, \Cref{lem:peeled-Q-count} &\raggedright \texttt{peel}, \texttt{peel\_\allowbreak correct} (\texttt{Section4/\allowbreak Peeled.lean});
\texttt{peelC\_\allowbreak multiplications}, \texttt{peelC\_\allowbreak additions},
\texttt{peelC\_\allowbreak multDepth} (\texttt{Height/\allowbreak PeeledCircuit.lean}) &
\raggedright Known monic powers and the binary recursion; exact counts
$2^{k-1}-1$ products, $5\cdot2^{k-2}-2$ additions, depth $k$ for $k\ge2$. &
full\\
\end{tabular}\par
\noindent\begin{tabular}{p{0.95in}p{2.05in}p{1.75in}p{0.55in}}
\raggedright \Cref{alg:constr-Tk2l} (structural layer) &\raggedright \texttt{Tpair}, \texttt{TF\_\allowbreak good}, \texttt{Tpair\_\allowbreak good}
(\texttt{Section5/\allowbreak T.lean}) &
\raggedright Monic/degree $k2^l$ for all four branches; no unit hypothesis. &
full\\
\end{tabular}\par
\noindent\begin{tabular}{p{0.95in}p{2.05in}p{1.75in}p{0.55in}}
\raggedright \Cref{lem:Rk2l}, \Cref{lem:Rk2l-leading-coeff},
\Cref{lem:Rk2l-top-boundary} &\raggedright \texttt{Rk2l\_\allowbreak triangular}, \texttt{Rk2l\_\allowbreak top\_\allowbreak two}
(\texttt{Section5/\allowbreak Rk2lTriMaster.lean}); \texttt{Rk2l\_\allowbreak deg},
\texttt{Rk2l\_\allowbreak lead} (\texttt{Section5/\allowbreak Rk2l.lean}) &
\raggedright \texttt{[Nontrivial A]}; $1,\dots,k$ units of $R$; $l\ge2$; tower of known
monic powers; the scalar-difference clause at $l=2$ for odd $k\ge3$. &
full\\
\end{tabular}\par
\noindent\begin{tabular}{p{0.95in}p{2.05in}p{1.75in}p{0.55in}}
\raggedright \Cref{lem:causal-perturbed-T} &\raggedright \texttt{causal\_\allowbreak perturbed\_\allowbreak T}, \texttt{Tpair\_\allowbreak compatiblePair}
(\texttt{Section5/\allowbreak PerturbedT.lean}) &
\raggedright $1,\dots,M$ units of $R$ ($M$ even); $l\ge2$. &
full\\
\end{tabular}\par
\noindent\begin{tabular}{p{0.95in}p{2.05in}p{1.75in}p{0.55in}}
\raggedright \Cref{lem:4k+1-splittable}, \Cref{lem:Q4k+1-from-H2} &\raggedright \texttt{fourk\_\allowbreak decodable}, \texttt{q4k1\_\allowbreak decodable}
(\texttt{Section6/\allowbreak GadgetDecoders.lean}); crown in
\texttt{Section5/\allowbreak FourKPlusOne.lean} &
\raggedright $1,\dots,2k$ units of $R$ and \texttt{IsUnit (2 : R)}; five explicit
$V$-relative pivots. &
full\\
\end{tabular}\par
\noindent\begin{tabular}{p{0.95in}p{2.05in}p{1.75in}p{0.55in}}
\raggedright \Cref{lem:Q-odd-degree-with-powers} &\raggedright \texttt{q\_\allowbreak odd\_\allowbreak degree\_\allowbreak decodable} (\texttt{Section6/\allowbreak QOddDegree.lean}) &
\raggedright $1,\dots,2k$ units of $R$; $l\ge2$; composes \texttt{fill\_\allowbreak correct},
\texttt{causal\_\allowbreak perturbed\_\allowbreak T}, \texttt{peel\_\allowbreak correct}. &
full\\
\end{tabular}\par
\noindent\begin{tabular}{p{0.95in}p{2.05in}p{1.75in}p{0.55in}}
\raggedright \Cref{lem:barQ15}, \Cref{lem:barQ8k+7} &\raggedright \texttt{barredGadgets\_\allowbreak of\_\allowbreak admissible} (\texttt{Examples/\allowbreak BarredGadgets.lean};
\texttt{BarQ15.lean}, \texttt{BarQGeneral.lean}) &
\raggedright \texttt{[Nontrivial A]}; $1,\dots,\mathrm{cap}$ units of $R$; the
$4\times4$ block with determinant $-k^2$ is inverted by adjugate. &
full\\
\end{tabular}\par
\noindent\begin{tabular}{p{0.95in}p{2.05in}p{1.75in}p{0.55in}}
\raggedright \Cref{lem:odd-gadgets-H2H4} &\raggedright \texttt{odd\_\allowbreak gadget\_\allowbreak dispatch} (\texttt{Section6/\allowbreak Dispatch.lean}) &
\raggedright $1,\dots,d$ units of $R$; hypothesis \texttt{BarredGadgets d}, discharged
by the previous row. &
full\\
\end{tabular}\par
\noindent\begin{tabular}{p{0.95in}p{2.05in}p{1.75in}p{0.55in}}
\raggedright \Cref{lem:8k+3-splittable}, \Cref{lem:8k+7-splittable} &\raggedright \texttt{eightk3\_\allowbreak compatible}, \texttt{eightk3\_\allowbreak decodable},
\texttt{eightk7\_\allowbreak compatible}, \texttt{eightk7\_\allowbreak decodable}
(\texttt{Section6/\allowbreak Induction.lean}) &
\raggedright \texttt{IsUnit (2 : R)} plus the smaller pair's compatibility and decoder
and the two gadget decoders as hypotheses. &
full\\
\end{tabular}\par
\noindent\begin{tabular}{p{0.95in}p{2.05in}p{1.75in}p{0.55in}}
\raggedright \Cref{lem:base-three-compatible} &\raggedright \texttt{base\_\allowbreak three\_\allowbreak compatible} (\texttt{Section6/\allowbreak SpecialCases.lean}) &
\raggedright None. &
full\\
\end{tabular}\par
\noindent\begin{tabular}{p{0.95in}p{2.05in}p{1.75in}p{0.55in}}
\raggedright \Cref{lem:septic-base} &\raggedright \texttt{septic\_\allowbreak good}, \texttt{septic\_\allowbreak decodable}
(\texttt{Examples/\allowbreak Septic.lean}); \texttt{Cost.\allowbreak SepticProgram.program}
(four products, ten additions) &
\raggedright \texttt{IsUnit (2 : R)}.  The unpublished impossibility of a three-product
joint septic realization (the reason $n=7$ is exceptional) has no Lean
counterpart. &
full\\
\end{tabular}\par
\noindent\begin{tabular}{p{0.95in}p{2.05in}p{1.75in}p{0.55in}}
\raggedright \Cref{lem:special-cases-splittable} ($15,27,31$) &\raggedright \texttt{decodable} in \texttt{Examples/\allowbreak P15.lean},
\texttt{Examples/\allowbreak P27Full.lean}, \texttt{Examples/\allowbreak P31Full.lean}
(namespaces \texttt{P15}, \texttt{P27Full}, \texttt{P31Full}) &
\raggedright \texttt{[Nontrivial A]}; unit hypotheses as consumed by the inner
\texttt{barQ15}, $Q_7$, $Q_3$ discharges. &
full\\
\end{tabular}\par
\noindent\begin{tabular}{p{0.95in}p{2.05in}p{1.75in}p{0.55in}}
\raggedright \Cref{lem:fill-Q-count}, \Cref{lem:T-multiplication-count},
\Cref{lem:odd-gadgets-count} &\raggedright \texttt{fill\_\allowbreak count}, \texttt{mers\_\allowbreak multiplication\_\allowbreak count},
\texttt{t\_\allowbreak multiplication\_\allowbreak count} (\texttt{Cost/\allowbreak Counts.lean});
\texttt{fourKPlusOne\_\allowbreak multiplication\_\allowbreak count},
\texttt{knownPowersOdd\_\allowbreak multiplication\_\allowbreak count},
\texttt{barredEightKPlusSeven\_\allowbreak multiplication\_\allowbreak count}
(\texttt{Cost/\allowbreak Gadgets.lean}) &
\raggedright Over $\mathbb N$ (schedule recurrences); the realized branch circuits
carry the same counts by construction. &
full\\
\bottomrule
\end{tabular}\par
\endgroup

\subsection{Lower bounds}
\label{sec:formalization-map:lower}

\begingroup
\footnotesize
\setlength{\tabcolsep}{3pt}
\par\noindent\begin{tabular}{p{0.95in}p{2.05in}p{1.75in}p{0.55in}}
\toprule
Paper & Lean (file) & Hypotheses in Lean & Status\\
\midrule
\end{tabular}\par\nopagebreak
\noindent\begin{tabular}{p{0.95in}p{2.05in}p{1.75in}p{0.55in}}
\raggedright Lemma ($6$ parameters), \Cref{sec:lower} &\raggedright \texttt{no\_\allowbreak rationalInverse\_\allowbreak general},
\texttt{no\_\allowbreak rationalInverse\_\allowbreak general\_\allowbreak of\_\allowbreak ringChar\_\allowbreak ne\_\allowbreak two}
(\texttt{LowerBound/\allowbreak General/\allowbreak Main.lean}); normal form
\texttt{exists\_\allowbreak singular\_\allowbreak jacobian},
\texttt{no\_\allowbreak rationalInverse\_\allowbreak affine}
(\texttt{LowerBound/\allowbreak Main.lean}); reduction steps
\texttt{gout\_\allowbreak eq} (\texttt{General/\allowbreak Gauge.lean}),
\texttt{polyJacobian\_\allowbreak outPolyGeneral} (\texttt{Affine.lean}),
\texttt{det\_\allowbreak eq\_\allowbreak zero\_\allowbreak of\_\allowbreak mulVec\_\allowbreak eq\_\allowbreak zero} (\texttt{DQ.lean}),
\texttt{det\_\allowbreak eq\_\allowbreak zero\_\allowbreak along\_\allowbreak orbit} (\texttt{Orbit.lean}),
\texttt{Qval\_\allowbreak Theta} (\texttt{Transversal.lean}); cross-checks
\texttt{det\_\allowbreak eq\_\allowbreak zero\_\allowbreak of\_\allowbreak Qval\_\allowbreak eq} (\texttt{Midpoint.lean}),
\texttt{RationalInverse.\allowbreak transport} (\texttt{Transport.lean}, outside the umbrella) &
\raggedright \texttt{[Field F] [Infinite F]}; \texttt{(2 : F) $\ne$ 0} (resp.\
\texttt{ringChar F $\ne$ 2}); six distinct points; the general three-gate
model \texttt{GCircuit F} (sixteen fixed constants, both first-gate constants
are slots, $A=B=0$ and $A=0\ne B$ allowed) and any affine map $Mp+h_0$ from
six parameters to seven slots, $M$ a $7\times6$ matrix with no rank condition.
The gauge identity, the chain rule for the quadratic slot map, cases
(i)--(iv) of \Cref{appendix:lower} (assembled by the characteristic-free pivot
route, so that $2\ne0$ enters only through the normal-form theorem) and the
first-gate degeneracies are proved; only an everywhere-defined rational left
inverse is excluded.  \texttt{GCircuit} directly encodes the topologically
ordered gate display of \Cref{appendix:lower}; no translation from an
arbitrary-program syntax is claimed. &
full\\
\end{tabular}\par
\noindent\begin{tabular}{p{0.95in}p{2.05in}p{1.75in}p{0.55in}}
\raggedright \Cref{thm:char2-lower} &\raggedright \texttt{no\_\allowbreak surjective\_\allowbreak eval}, \texttt{no\_\allowbreak construction\_\allowbreak general}
(\texttt{LowerBoundChar2/\allowbreak General.lean}); affine case
\texttt{no\_\allowbreak construction} (\texttt{LowerBoundChar2/\allowbreak Main.lean}); sharpness
\texttt{oneGate\_\allowbreak isConstruction} (\texttt{LowerBoundChar2/\allowbreak Sharpness.lean}) &
\raggedright \texttt{[Field F] [Fintype F] [CharP F 2]}; $1<n$; $2n\le|F|$; an $n$-gate
chain \texttt{Circuit F n}; an arbitrary preprocessing map
\texttt{pre : $\Lambda$ $\to$ Slots F n}.  The fibre-counting proof of
\Cref{sec:lower-char2} is formalized as stated; the affine case is a
separate, earlier development. &
full\\
\end{tabular}\par
\noindent\begin{tabular}{p{0.95in}p{2.05in}p{1.75in}p{0.55in}}
\raggedright \Cref{lem:first-char2-circuit-inverse} &\raggedright \texttt{circuit\_\allowbreak eval\_\allowbreak bijective} (\texttt{Examples/\allowbreak Char2Inverse.lean}) &
\raggedright \texttt{[Field F] [CharP F 2] [PerfectRing F 2]}; seven distinct points. &
full\\
\end{tabular}\par
\noindent\begin{tabular}{p{0.95in}p{2.05in}p{1.75in}p{0.55in}}
\raggedright \Cref{lem:char2-small-staircase-butterfly} (degree $9$) &\raggedright
\texttt{coeffMap9Fin\_\allowbreak bijective},
\texttt{P9\_\allowbreak eval\_\allowbreak bijective}
(\texttt{Examples/\allowbreak Char2SmallInverses.lean}) &
\raggedright \texttt{[Field F] [CharP F 2]}; nine distinct points.  The
literal circuit identity, both directions of the displayed coefficient
decoder, and the Vandermonde evaluation step are checked.  The degree-$11$
half of the paper lemma is not yet formalized. &
partial\\
\bottomrule
\end{tabular}\par
\endgroup

\subsection{Unit hypotheses versus \texorpdfstring{$n$}{n}-admissibility}
\label{sec:formalization-map:units}

The paper's standing hypothesis is that $\F$ is $n$-admissible
(\Cref{appendix:constructions}: characteristic zero or characteristic
$p>n$), which \Cref{rem:char2} describes as sufficient rather than sharp.
The Lean master theorem \texttt{odd\_\allowbreak realizable\_\allowbreak pairs} assumes instead,
over an arbitrary commutative ring $R$ (no field, no characteristic),
\begin{center}
  \texttt{$\forall$ i : $\mathbb N$, 1 $\le$ i $\to$ i $\le$ n $\to$ IsUnit (((i : $\mathbb N$) : $\mathbb Z$) : R)},
\end{center}
that is, exactly the integers $1,2,\dots,n$ are units of $R$; this is the
predicate \texttt{Admissible R n} of \texttt{Admissible.lean} up to the
integer cast (\texttt{Int.cast\_\allowbreak natCast}).  Every composite pivot of the
construction, namely the scalar slopes dividing $2k(k-1)$, the block
determinant $-k^2$ of \Cref{lem:barQ8k+7}, and the factor $2$ of every
square gadget, is inverted in Lean as a product of unit factors that are
each at most $n$ (\texttt{Admissible.isUnit\_\allowbreak mul\_\allowbreak cast}); no hypothesis on
an integer larger than $n$ is ever used.  The sub-lemmas consume the same
predicate restricted to their own degree (for instance $1,\dots,k$ in
\texttt{Rk2l\_\allowbreak triangular}, $1,\dots,2k$ in \texttt{fourk\_\allowbreak decodable},
$1,\dots,d$ in \texttt{odd\_\allowbreak gadget\_\allowbreak dispatch}), and several need only
\texttt{IsUnit (2 : R)} or a single root index, as the tables record.  For a
field $\F$ the Lean hypothesis is equivalent to $n$-admissibility: it holds
in characteristic zero and in characteristic $p>n$
(\texttt{admissible\_\allowbreak of\_\allowbreak charZero}, \texttt{admissible\_\allowbreak of\_\allowbreak charP}), and it
fails in characteristic $p\le n$ because $p$ itself is then not a unit.
Thus the Lean statements do not certify the finer claim of \Cref{rem:char2}
that only the pivots actually displayed need be nonzero; they certify the
uniform sufficient condition.  The coverage endpoints
(\texttt{monic\_\allowbreak coefficient\_\allowbreak map\_\allowbreak bijective} and the automorphism of
\Cref{lem:polynomial-left-inverse-automorphism}) additionally assume the
base ring is Noetherian and nontrivial, which every field satisfies.  The
lower bounds use the field-level hypotheses shown in their table: an
infinite field with $2\ne0$ for degree six, and a finite field of
characteristic two for \Cref{thm:char2-lower}.

\paragraph{Build and axioms.}
\begin{sloppypar}
The construction, cost, height, example, and both lower-bound
developments (including \texttt{HeightFinal.lean},
\texttt{LowerBound/\allowbreak General/\allowbreak Main.lean},
\texttt{LowerBoundChar2/\allowbreak General.lean}, and
\texttt{LowerBoundChar2/\allowbreak Sharpness.lean}) are imported by the umbrella
module \texttt{FastPoly.lean}, the default target; \texttt{lake build FastPoly}
completes it with no errors. No file under \texttt{FastPoly/\allowbreak } declares an \texttt{axiom} or
contains a \texttt{sorry}.  \texttt{\#print axioms} on
\texttt{polynomial\_\allowbreak height}, \texttt{polynomial\_\allowbreak height\_\allowbreak of\_\allowbreak charZero},
\texttt{polynomial\_\allowbreak height\_\allowbreak of\_\allowbreak charP} (\texttt{HeightFinal.lean}),
\texttt{decodedAdditionPolynomial\_\allowbreak exists}
(\texttt{Cost/\allowbreak Additions/\allowbreak DecodedPolynomial.lean}), and
\texttt{MainArrangementsChecked.of\_\allowbreak charZero},
\texttt{MainArrangementsChecked.of\_\allowbreak charP} (\texttt{PaperMain.lean}), as well as on
\texttt{no\_\allowbreak surjective\_\allowbreak eval}, \texttt{no\_\allowbreak construction\_\allowbreak general},
\texttt{no\_\allowbreak construction'} (\texttt{LowerBoundChar2/\allowbreak General.lean}),
\texttt{no\_\allowbreak rationalInverse\_\allowbreak general}
(\texttt{LowerBound/\allowbreak General/\allowbreak Main.lean}), and
\texttt{oneGate\_\allowbreak isConstruction}
(\texttt{LowerBoundChar2/\allowbreak Sharpness.lean})
reports only \texttt{propext}, \texttt{Classical.choice}, and
\texttt{Quot.sound}; the same is recorded for
\texttt{construction\_\allowbreak additions\_\allowbreak checked}
(\texttt{Cost/\allowbreak Additions/\allowbreak Realization.lean};
\texttt{FastPoly/\allowbreak ROADMAP.md}) and for the affine characteristic-two
development (module docstring of \texttt{LowerBoundChar2/\allowbreak Main.lean}).
No separate \texttt{\#print axioms} transcript is recorded for the remaining
theorems.
\end{sloppypar}
